\documentclass[10pt,a4paper]{amsart}
\usepackage[margin=19mm]{geometry}
\usepackage{amsmath,amssymb,mathtools}
\usepackage[scr=rsfs,cal=euler]{mathalfa}
\usepackage{xfrac}
\usepackage[unicode,hidelinks,bookmarks=false]{hyperref}
\newcommand{\ie}{i.e.\ }
\newcommand{\cf}{cf.\ }
\newcommand{\resp}{respectively\ }
\newcommand{\Subsection}{\subsection}
\providecommand{\nobreakdash}{\nobreak\hbox{-}}
\setlength{\emergencystretch}{2em}
\multlinegap0pt
\usepackage[all]{xy}
\usepackage{amssymb}
\usepackage{xcolor}
\newtheorem{thm}{Theorem}
\newtheorem{prop}{Proposition}
\DeclareMathOperator{\GL}{GL}
\DeclareMathOperator{\Sh}{Sh}
\def\bC{\mathbb{C}}
\def\cL{\mathcal{L}}
\def\cM{\mathcal{M}}
\DeclareMathOperator{\rG}{G}
\DeclareMathOperator{\rH}{H}
\DeclareMathOperator{\rN}{N}
\DeclareMathOperator{\sh}{sh}
\title{Hyperbolicity in presence of a large local system}
\author{Yohan Brunebarbe}
\date{Source: arXiv:2207.03283v2 (CC BY 4.0)}
\begin{document}
\maketitle

\noindent\emph{Source and licence.} Hyperbolicity in presence of a large local system. Version arXiv:2207.03283v2 (CC BY 4.0), distributed by arXiv under the Creative Commons Attribution 4.0 International licence, http://creativecommons.org/licenses/by/4.0/, as recorded in the arXiv licence metadata for this exact version and shown on the abstract page https://arxiv.org/abs/2207.03283v2. Full licence notice: \texttt{licenses/arxiv-2207.03283-CC-BY-4.0.txt}.\par\smallskip

\noindent\textit{Editorial scope.} Counted unit: the article's prop solvable/semisimple (source0.tex lines 397--404). The article's complete original proof of it is delivered with it (source0.tex lines 405--411, one \texttt{proof} environment of 7 lines). No mathematics is written, repaired or re-derived: the statement and the proof are the article's own lines, in the article's own notation, and no retained line is substituted or re-flowed.  Retained uncounted as the source-local context the proof needs: lines 312--317; lines 321--329; lines 365--367; lines 381--384.  Nothing was omitted from any retained span, so the package carries no omission record.  No retained line contains a citation, so the wrapper carries no bibliography and no bibliography entry.  No retained line contains a figure environment, an image-inclusion command or drawing code, so no image is delivered and the package declares no asset dependency.  Editorial formatting only, in addition: an AMS \texttt{amsart} preamble that restates the article's own declarations and macros used by the retained lines, namely the environments \texttt{thm}, \texttt{prop} on separate counters, together with the standard packages those lines need.  The retained environments print in this document as Theorem 1, Proposition 1 in order of appearance, whereas the article prints the same items with the numbering of its own full document.  The complete native source of the article as distributed in the arXiv e-print for this exact version is bundled unchanged as \texttt{sources/127907/source0.tex}.
\begin{thm}\label{existence of Shafarevich morphism}
Let $\cL$ be a complex local system on a proper irreducible normal complex algebraic variety $X$. Assume that $\cL$ is semisimple, or equivalently that the algebraic monodromy group of $\cL$ is reductive.\\
Then, there exist a projective normal complex algebraic variety $\Sh_X^\cL$ and a fibration $\sh_X^\cL \colon X \rightarrow \Sh_X^\cL$, unique up to a unique isomorphism, with the following property: for any connected proper complex algebraic variety $Z$ and any morphism $f \colon Z \rightarrow X$, the composite map $\sh_X^\cL \circ f \colon Z \rightarrow \Sh_X^\cL$ is constant if and only if the local system
$f^{-1} \cL$ is isotrivial. \\
Moreover, if the monodromy group of $\cL$ is torsion-free, then there exists a (necessarily large) complex local system $\cM$ on $\Sh_X^\cL$ such that $\cL = \left(\sh_X^\cL \right)^{-1} \cM$.
\end{thm}

\begin{prop}\label{Shafarevich: from smooth to normal}
Let $X^\prime \to X$ be a fibration between normal proper complex algebraic varieties. Let $\cL$ be a semisimple complex local system on $X$ and $\sh_{X^\prime}^{\cL^\prime} \colon X^\prime \rightarrow \Sh_{X^\prime}^{\cL^\prime}$ be the Shafarevich morphism associated to the (necessarily semisimple) local system $ \cL^\prime := \nu^{-1} \cL$. Then there is a (unique) factorization
\begin{align*}
\xymatrix{
X^\prime \ar[r] \ar[d] &  \Sh_{X^\prime}^{\cL^\prime} \\
X   \ar[ru]       &   }
\end{align*}
and the induced morphism $X \rightarrow \Sh_{X^\prime}^{\cL^\prime}$ is the Shafarevich morphism associated to $\cL$.
\end{prop}

\begin{prop}\label{Shafarevich in the solvable case}
Let $\cL$ be a large complex local system on a normal projective complex algebraic variety $X$. Let $\pi_1(X) \to \Gamma \subset \GL(n , \bC)$ be its monodromy representation. Assume that its monodromy group $\Gamma$ is solvable and the commutator subgroup $[\Gamma, \Gamma]$ is nilpotent. Then, $X$ admits a finite morphism to an abelian variety.
\end{prop}

\begin{prop}\label{criterion derived group nilpotent}
If $\Gamma$ is a subgroup of $ \GL(n , \bC)$ contained in a connected solvable algebraic subgroup of $ \GL(n , \bC)$, then
$[\Gamma, \Gamma]$ is nilpotent.
\end{prop}

\begin{prop}\label{fibration solvable/semisimple}
Let $X$ be a normal irreducible proper complex algebraic variety supporting a complex local system $\cL$. Then, there exist a finite étale cover $X^\prime \to X$ and a fibration $f \colon X^\prime \to Y$ onto a normal irreducible projective complex algebraic variety $Y$ such that:
\begin{itemize}
\item $Y$ admits a large complex local system with torsion-free monodromy and a semisimple algebraic monodromy group;
\item the monodromy of the restriction of $\cL$ to any fiber of $f$ is solvable.
\end{itemize} 
Moreover, if $\cL$ is large (resp. big), then the normalization of every fiber of $f$ admits a finite morphism to an abelian variety (resp. every generic fiber of $f$ admits a morphism to an abelian variety which is generically finite onto its image).
\end{prop}

\begin{proof}
Consider the monodromy representation $\pi_1(X) \to \rG(\bC)$ of $\cL$,  where $\rG$ is the algebraic monodromy group of $\cL$. Up to replacing $X$ with a finite étale cover, one can assume that $\rG$ is connected. Let $\rN$ be the solvable radical of $\rG$, so that $\rN$ is a connected solvable complex algebraic group. Let $\rH$ be the quotient of $\rG$ by $\rN$, so that $\rH$ is a connected semisimple complex algebraic group. The induced representation $\pi_1(X) \to \rH(\bC)$ has a Zariski-dense image, and replacing $X$ with a finite étale cover, one can assume that it has torsion-free image thanks to Selberg lemma. \\
Let $f \colon X \to Y$ denote the Shafarevich morphism associated to $\pi_1(X) \to \rH(\bC)$, cf. Theorem \ref{existence of Shafarevich morphism}. In particular, $Y$ is a normal irreducible projective complex algebraic variety and $f$ is surjective with connected fibres.
Since the representation $\pi_1(X) \to \rH(\bC)$ has torsion-free image, it factors through the homomorphism $\pi_1(X) \to \pi_1(Y)$, cf. Theorem \ref{existence of Shafarevich morphism}. The induced homomorphism $\pi_1(Y) \to \rH(\bC)$ corresponds to a large complex local system with a semisimple algebraic monodromy group.\\
Let $F$ be a (necessarily connected) fiber of $f$. Since the induced morphism $F \to Y$ is constant, the composite homomorphism $\pi_1(F) \to \pi_1(X) \to \rH(\bC)$ has finite image. But the image of $\pi_1(X) \to \rH(\bC)$ is torsion-free, hence the image of $\pi_1(F) \to \rH(\bC)$ is in fact trivial. Therefore it is contained in $\rN(\bC)$, from what it follows that it is solvable.\\
Assume now in addition that $\cL$ is large, and let $F^\prime$ be the normalization of an irreducible component of a fiber of $f$. Since the monodromy of the restriction of $\cL$ to $F^\prime$ is contained in the connected solvable algebraic group $\rN(\bC)$, one can apply Proposition \ref{Shafarevich: from smooth to normal}, Proposition \ref{Shafarevich in the solvable case} and Proposition \ref{criterion derived group nilpotent} to infer that $F^\prime$ admits a finite morphism to an abelian variety. The statement follows by taking a product over the irreducible components. The case where $\cL$ is big is similar and left to the reader.
\end{proof}

\end{document}
